% !TeX program = pdflatex
% Overleaf-ready. Klasse bei Bedarf gegen Konferenz-/Zeitschriftenvorlage tauschen.
\documentclass[11pt,a4paper]{article}
\usepackage[T1]{fontenc}
\usepackage[utf8]{inputenc}
\usepackage[ngerman]{babel}
\usepackage{csquotes}
\usepackage[style=authoryear,backend=biber]{biblatex}
\usepackage{booktabs}
\usepackage{graphicx}
\usepackage{hyperref}
\addbibresource{refs.bib}

\title{MLLM-basierte Informationsextraktion aus historischen Kirchenbüchern:\\
Eine Evaluation von Gemini~3.5~Flash gegen die Matrikel-Datenbank GenDB}
\author{Tobias Perschl\\Universität Passau}
\date{\today}

\begin{document}
\maketitle

\begin{abstract}
% TODO: 150--250 Wörter. Kernaussagen: (1) Verkartungsdatenbanken wie GenDB
% lassen sich als Gold Standard für die Evaluation multimodaler LLMs
% zweckentfremden; (2) wir evaluieren Gemini 3.5 Flash auf der Extraktion
% strukturierter Personendaten direkt aus Matrikel-Scans (Taufen, Heiraten,
% Sterbefälle, 17.--19. Jh., Kurrent/Latein); (3) zentrale Ergebnisse
% (Entry-Detection-F1, feldweise Accuracy/CER, Ablationen).
\end{abstract}

\section{Einleitung}
% TODO: Motivation. Kirchenbücher als zentrale Quelle für Genealogie und
% historische Demographie; manuelle Verkartung ist teuer; MLLMs versprechen
% End-to-End-Extraktion ohne separates HTR/NER-Pipelining. Forschungsfrage:
% Wie gut extrahiert ein multimodales LLM die Felder einer professionellen
% Verkartung direkt aus dem Scan -- und was sagt der Vergleich über beide aus?
% Beiträge: (1) reproduzierbare Korpus-Pipeline GenDB/Matricula,
% (2) Evaluationsrahmen mit Entry-Alignment und feldsensitiver Bewertung,
% (3) empirische Ergebnisse inkl. Ablationen, (4) methodische Reflexion
% "Index vs. Quelle".

\section{Forschungsstand}
\subsection{HTR und Dokumentanalyse historischer Handschriften}
% TODO: klassische Pipelines: Layoutanalyse + HTR + nachgelagerte IE.
% Transkribus \parencite{muehlberger2019transforming}, TrOCR
% \parencite{li2023trocr}, OCR-freie End-to-End-Modelle wie Donut
% \parencite{kim2022donut}.
\subsection{Multimodale LLMs für historische Dokumente}
% TODO: jüngere Arbeiten zu LLM-/MLLM-Transkription historischer
% Handschriften, z.\,B. \parencite{humphries2024unlocking}; Abgrenzung:
% wir evaluieren nicht Transkription, sondern direkte schemabasierte
% Informationsextraktion (Structured Output) gegen eine unabhängig
% entstandene Verkartung.
\subsection{Kirchenbuchdatenbanken als Forschungsdaten}
% TODO: Matricula Online, GenDB (Bistumsarchiv Passau / FORWISS);
% Verkartung vs. Transkription; bekannte Eigenschaften solcher Indizes
% (Normalisierung, Schätzdaten, selektive Felder).

\section{Daten}
\subsection{GenDB und Matricula}
% TODO: Beschreibung der Datenbank (Eintragstypen T/M/S, Feldschema),
% Verknüpfung zu Matricula-Scans.
\subsection{Korpuskonstruktion}
Für jede ausgewählte Registerseite werden über die GenDB-Suchschnittstelle
sämtliche verkarteten Einträge aller drei Eintragstypen abgefragt, mit dem
zugehörigen Matricula-Scan verknüpft und in einem maschinenlesbaren Manifest
zusammengeführt (ein Bild $\rightarrow$ $n$ Gold-Einträge).
% TODO: Stichprobenziehung (Pfarreien, Zeiträume, Eintragstypen,
% Tabellen- vs. Fließtextlayout), Korpusstatistik als Tabelle.
\subsection{Eigenschaften des Gold Standards}
Die Verkartung ist ein \emph{Index}, keine diplomatische Transkription: Felder
können fehlen, obwohl die Information im Scan vorhanden ist (z.\,B. nur ein
geschätztes Jahr trotz tagesgenauer Datumsangabe in der Quelle), und
Schreibweisen können normalisiert sein. Die Evaluation bewertet daher
ausschließlich Felder, die im Gold Standard belegt sind.

\section{Methode}
\subsection{Extraktion}
Gemini~3.5~Flash erhält den Seitenscan und extrahiert Einträge per Structured
Output gegen ein JSON-Schema, das die GenDB-Felder spiegelt. Zwei Settings:
\emph{Page-Level} (alle Einträge einer Seite in einem Aufruf) und
\emph{Entry-Level} (ein Eintrag pro Aufruf, lokalisiert über laufende Nummer
bzw. Position). Experimentvariablen: \texttt{thinking\_level},
\texttt{media\_resolution}, Few-Shot-Beispiele.
\subsection{Alignment und Metriken}
Im Page-Setting werden vorhergesagte und Gold-Einträge je Seite über eine
gewichtete Ähnlichkeit (Namen, Brautnamen, Datum, Nummer, Eintragstyp;
normalisiert über die im Gold belegten Anker) optimal 1:1 zugeordnet
(ungarische Methode). Daraus ergeben sich Entry-Detection-Precision/-Recall/-F1
(ungematchte Gold-Einträge = Misses, ungematchte Vorhersagen =
Halluzinationen). Auf gematchten Paaren werden Felder bewertet: Namens- und
Textfelder mit normalisiertem Exact Match und Character Error Rate, numerische
und kategoriale Felder mit Accuracy; Datumsfelder zusätzlich getrennt für
nicht-geschätzte Gold-Daten.

\section{Ergebnisse}
% TODO: Haupttabelle (Entry-Detection + feldweise Acc/CER, Page vs. Entry),
% Ablationen (thinking_level, media_resolution, Few-Shot), Aufschlüsselung
% nach Eintragstyp und Jahrhundert, qualitative Fehleranalyse mit Beispielen
% aus comparisons_*.csv.
\begin{table}[ht]
\centering
\caption{Feldweise Ergebnisse (Platzhalter).}
\begin{tabular}{lrrr}
\toprule
Feld & $n$ & Accuracy & CER \\
\midrule
Nachname & -- & -- & -- \\
Vorname & -- & -- & -- \\
Jahr & -- & -- & -- \\
\bottomrule
\end{tabular}
\end{table}

\section{Diskussion}
% TODO: (1) Wo scheitert das MLLM (Kurrent-Lesung, Spaltenzuordnung,
% Eintragssegmentierung, Halluzinationen)? (2) Wo ist die Vorhersage
% plausibel korrekt, aber der Index lückenhaft -- Umkehrung der
% Gold-Standard-Annahme; Potenzial des MLLM zur Index-Anreicherung.
% (3) Normalisierung vs. originalgetreue Schreibweise. (4) Kosten/Nutzen
% gegenüber klassischen HTR-Pipelines. (5) Limitationen: Stichprobe,
% nur ein Modell, Verkartungsqualität.

\section{Fazit und Ausblick}
% TODO.

\section*{Daten- und Codeverfügbarkeit}
Die vollständige Pipeline (Scraper, Extraktion, Evaluation) ist verfügbar
unter \url{https://github.com/Maelkolb/MLLM_IE_Kirchenmatrikel}. % TODO: Repo-URL prüfen
Bildrechte der Scans liegen bei Matricula Online bzw. den Archiven; das
Korpus wird als Manifest (Seitenreferenzen + Verkartungsdaten) publiziert.

\printbibliography
\end{document}
